\chapter{Il teacher: un sistema a regole}\label{cap:teacher}

% Capitolo 4 — Il teacher: un sistema a regole.

Prima di parlare di reti neurali e training, parliamo di chi "insegna": il
\textbf{teacher rule-based}. È il componente che assegna la classe di rischio
a ogni combinazione di parametri vitali del dataset sintetico (capitolo
\ref{cap:dati}). L'MLP, nei capitoli successivi, impara ad approssimare
proprio questa funzione: è il punto di partenza della knowledge distillation
(capitolo \ref{cap:distillazione}).

\section{Cos'è un sistema a regole}\label{sec:cos-e-sistema-a-regole}

Un sistema a regole è il modo più semplice e antico di scrivere un
"classificatore": una serie di \codice{if}/\codice{elif} che confrontano i
valori in ingresso con soglie e range predefiniti e restituiscono una
decisione. È il tipo di codice che un medico potrebbe scrivere a mano,
traducendo le linee guida cliniche in condizioni esplicite.

\begin{lstlisting}[caption={La struttura di un sistema a regole (pseudocodice)}]
if valore >= soglia_critica:
    classe = "alto"
elif valore >= soglia_moderata:
    classe = "medio"
else:
    classe = "basso"
\end{lstlisting}

Un sistema a regole non impara nulla: ogni decisione è scritta esplicitamente
nel codice. Questo lo rende deterministico (stesso input, stessa risposta,
sempre) e completamente trasparente (ogni decisione è riconducibile a una
condizione leggibile).

\section{Le regole cliniche del progetto}\label{sec:regole-cliniche-progetto}

Nel progetto le regole cliniche vivono in un unico posto:
\file{src/telemedicina\_supervised/safety/safety\_rules.py}. È la
\textbf{fonte unica di verità} delle soglie: il teacher offline, il
validatore e il safety gate runtime dipendono tutti da questo modulo, e
nessuna soglia numerica è duplicata altrove (un guard test,
\file{tests/test\_single\_source.py}, lo verifica automaticamente).

La logica di classificazione è la seguente:

\begin{enumerate}
  \item \textbf{Validità}: se un valore è fuori dal range fisiologico, o se
        la sistolica non è maggiore della diastolica, la risposta è
        \classe{errore}.
  \item \textbf{Anomalie}: per ogni parametro fuori dal range normale si
        calcola la gravità: \emph{critica} (oltre le soglie critiche),
        \emph{moderata} (oltre le soglie moderate) o \emph{lieve}.
  \item \textbf{Pattern critici}: alcune \emph{combinazioni} di parametri
        sono pericolose anche se ogni singolo valore non è critico (es.
        ipotensione con tachicardia = shock).
  \item \textbf{Decisione}: se c'è un pattern critico o almeno un'anomalia
        critica, la classe è \classe{alto}; se c'è qualunque anomalia non
        critica, \classe{medio}; altrimenti \classe{basso}.
\end{enumerate}

Le soglie critiche reali del progetto sono:

\begin{table}[h]
\centering
\caption{Soglie critiche per parametro (da \file{safety\_rules.py}).}
\label{tab:soglie-critiche}
\begin{tabular}{ll}
\toprule
\textbf{Parametro} & \textbf{Soglia critica} \\
\midrule
\codice{pressione\_sistolica} & $\geq$ 180 mmHg \\
\codice{pressione\_diastolica} & $\geq$ 110 mmHg \\
\codice{frequenza\_cardiaca} & $\geq$ 120 bpm oppure $\leq$ 50 bpm \\
\codice{temperatura} & $\geq$ 38.5 °C oppure $\leq$ 35.0 °C \\
\codice{saturazione\_ossigeno} & $\leq$ 90 \% \\
\codice{glicemia} & $\geq$ 200 mg/dL oppure $\leq$ 60 mg/dL \\
\bottomrule
\end{tabular}
\end{table}

Il teacher vero e proprio è in \file{training/teacher\_rules.py}: la classe
\codice{TeacherRules} con il metodo \codice{label}. Il punto architetturale
centrale è che il teacher \emph{non possiede soglie proprie}: riusa la stessa
funzione \codice{classifica\_rischio} del modulo di sicurezza. In questo modo
le etichette (offline) e il safety gate (runtime) non possono divergere.

\begin{lstlisting}[caption={Il metodo label del teacher (codice reale, abbreviato)}]
# training/teacher_rules.py
from telemedicina_supervised.safety.safety_rules import classifica_rischio

class TeacherRules:
    """Teacher rule-based offline: produce la label per una
    combinazione di parametri vitali."""

    def label(self, parametri):
        """Label di rischio: 'basso' | 'medio' | 'alto' | 'errore'."""
        return classifica_rischio(parametri)

    # Alias esplicito per uso in generazione dataset.
    etichetta = label
\end{lstlisting}

\begin{esempio}{Un caso critico reale}
Con i parametri (185, 110, 80, 36.8, 98, 95) — sistolica 185 e diastolica
110 — le regole rilevano una crisi ipertensiva (pattern critico) e
rispondono \classe{alto}. Il safety gate runtime applica la stessa logica:
un caso del genere non passa mai dall'MLP.
\end{esempio}

\section{Vantaggi}\label{sec:vantaggi}

Il teacher rule-based ha quattro vantaggi che lo rendono ideale come fonte
delle etichette:

\begin{itemize}
  \item \textbf{Trasparenza totale}: ogni decisione è una condizione
        leggibile nel codice. Non c'è nessuna "scatola nera".
  \item \textbf{Determinismo}: stesso input, stessa risposta, sempre. Le
        etichette sono riproducibili per costruzione.
  \item \textbf{Spiegabilità}: per ogni caso si può dire \emph{perché} la
        classe è quella: basta elencare le anomalie e i pattern rilevati.
  \item \textbf{Nessun training}: non servono dati storici né un processo di
        apprendimento: le regole sono scritte a mano e funzionano subito.
\end{itemize}

\section{Limiti}\label{sec:limiti}

I sistemi a regole hanno però limiti ben noti, ed è importante dichiararli:

\begin{itemize}
  \item \textbf{Rigidità}: le regole coprono solo i casi previsti a priori.
        Una combinazione di parametri mai immaginata non viene gestita.
  \item \textbf{Soglie arbitrarie}: i confini tra \classe{basso},
        \classe{medio} e \classe{alto} sono numeri fissi (es. 140 mmHg per
        la sistolica), ma la realtà clinica è sfumata: un valore di 141 non
        è qualitativamente diverso da 139.
  \item \textbf{Manutenzione manuale}: ogni nuova conoscenza clinica va
        tradotta a mano in nuove condizioni, con il rischio di errori e
        incoerenze.
  \item \textbf{Confini netti}: la classificazione è a gradino: basta un
        decimo di grado per cambiare classe. È proprio per questo che il
        generatore del dataset usa la buffer zone (capitolo \ref{cap:dati}).
\end{itemize}

Questi limiti non rendono il teacher inutile: lo rendono il \emph{punto di
partenza} giusto. Le regole sono la conoscenza di base; il modello impara ad
approssimarle in modo fluido e continuo, e il safety gate resta a garantire
che nessun caso critico venga mai declassato.

\section{Il teacher come oracolo}\label{sec:teacher-oracolo}

Nel progetto il teacher ha un ruolo preciso: è l'\textbf{oracolo} che
etichetta il dataset sintetico. Ogni campione generato (capitolo
\ref{cap:dati}) viene passato a \codice{label}, e la classe restituita
diventa la label del campione. In questo modo il dataset contiene, per
costruzione, la "verità" del teacher: l'MLP non impara da dati clinici reali,
ma dalla funzione rule-based.

Questo è esattamente il punto di partenza della \textbf{knowledge
distillation} \cite{hinton2015distilling}: un sistema complesso (qui, un
sistema a regole) "dista" la propria conoscenza a un modello più semplice
(qui, un MLP). Il capitolo \ref{cap:distillazione} approfondisce questo
passaggio.

\begin{esempio}{La baseline del teacher}
Il report di training (\file{data/models/report.json}) registra la
\emph{baseline rule-based}: il teacher applicato alle proprie label sul test
ha un accordo di 1.0 su 20\,000 campioni. È l'upper bound di imitazione:
l'accuracy dell'MLP va letta come fedeltà alle regole, non come accuratezza
diagnostica su pazienti reali.
\end{esempio}